\section{Discussion}
\label{sec:discussion}

\subsection{Key Findings}

\textbf{Two mechanistically distinct pathways to WGA improvement.} DFR achieves WGA = 0.91 by recalibrating the head while leaving backbone features --- with their full spurious background encoding (ERM probe accuracy $\approx 0.984$) --- entirely unchanged. GroupDRO achieves WGA = 0.88 by modifying backbone representations such that background information is significantly less linearly extractable ($p = 0.0039$, $d = 6.48$). These are two distinct robustification pathways, not two points on a continuum.

This finding has a non-obvious implication: the backbone's spurious encoding level is neither necessary nor sufficient for WGA improvement in isolation. DFR demonstrates that high spurious encoding in the backbone does not prevent high WGA if the head is correctly calibrated. GroupDRO demonstrates that reducing backbone spurious encoding also yields high WGA. The mechanisms are alternative solutions to the same problem.

\textbf{GroupDRO backbone modification is substantial and consistent.} The backbone/head L2 ratio of 6.47--7.03 indicates that GroupDRO's group-reweighted gradient signal primarily reshapes backbone weights. The DFR negative control (backbone/head ratio = 0.000 for all seeds) eliminates checkpoint-specific artifact confounds.

\textbf{Connection to the literature.} Our finding that GroupDRO implicitly reduces background linear decodability aligns with Park et al.\ \cite{park2025spurious} (SCER explicitly regularizes spurious feature directions). We partially nuance Izmailov et al.\ \cite{izmailov2022feature}, who suggested WGA improvement is ``primarily in the head'' --- our H-M2 results show the backbone changes substantially (ratio 6.47--7.03), though this does not contradict the head's importance for WGA per se. Our weight-difference analysis corroborates Raymond et al.\ \cite{raymond2026groupdro}; our finding stands on its own pre-registered evidence independent of this citation.

\subsection{Limitations}

\textbf{L1: $n=3$ seeds.} One-sided t-test with $n=3$ has $\sim$55\% power for $d=0.8$. The primary finding ($d=6.48$) is unaffected. H-M2's linear probe result ($p=0.0526$) is borderline SUGGESTIVE, likely due to smaller probe training set than H-M3's full $N=5{,}794$.

\textbf{L2: Single dataset and architecture.} All experiments use Waterbirds WILDS with ResNet-50. Generalization to CelebA, MultiNLI, ViT, or DINO architectures is an open question.

\textbf{L3: H-P2 SUGGESTIVE at $n=9$.} The pre-registered CONFIRMED threshold (CI upper bound $< 0$) is not met. Root cause: method-level WGA constants create within-method collinearity inflating bootstrap CI width. The ERM+GroupDRO ablation is CONFIRMED ($r=-0.755$, $p=0.041$).

\textbf{L4: Decodability $\neq$ mechanism.} Reduced probe accuracy establishes that background information is less linearly extractable from GroupDRO features; it does not establish \emph{why}. ``Suppression vs.\ dilution'' mechanistic ambiguity cannot be resolved from probe accuracy alone.

\textbf{L5: Core attribute probe not measured.} We did not measure bird species (core attribute) probe accuracy from GroupDRO features. Whether GroupDRO preserves core-feature encoding is an important future validation.

Establishing causality would require intervention studies (e.g., ablating minority upweighting while holding architecture fixed); our results establish consistent mechanistic co-occurrence, not causal necessity.

\subsection{Broader Impact}

Our diagnostic framework --- spurious attribute linear probe accuracy on frozen backbone features --- is low-cost (no new training, forward-pass only), reproducible (standard sklearn tools), and directly interpretable. Applied systematically, it provides a backbone spurious encoding ``fingerprint'' for any robustification method with publicly available checkpoints.

Understanding the backbone-vs-head distinction informs method selection: when group labels for head retraining are available, DFR's simpler intervention may suffice. When backbone-level change is desired for transfer learning or when group annotations are unavailable at inference time, GroupDRO's backbone modification offers a substantively different form of robustification.
